\documentclass{situationist-preprint}

\title{Absolute Freedom and Facticity}
\msid{PHE-10-a}
\msfamily{PHE-10}
\msdiscipline{Phenomenology}
\mstype{Principal manuscript}
\mskeywords{Sartre; freedom; facticity; situation; coefficient of adversity; Merleau-Ponty; sedimentation; Beauvoir; project}

\begin{document}
\maketitle

\begin{abstract}
Sartre holds that human freedom is absolute: we are condemned to be free, and nothing outside our freedom can limit it. He also holds that freedom is always exercised in a situation, against a given that it did not create: a place, a past, an environment, other people and death. Critics, Merleau-Ponty first among them, have argued that the two claims are in tension, and that if freedom is everywhere it is nowhere. We argue that freedom is ontologically absolute but always exercised within facticity, and that the two are not in tension. The key move is to distinguish two questions: whether the facts of a situation determine what they mean for the agent, and whether they determine what the agent can bring about. Sartre's absoluteness is a negative answer to the first question; facticity is a positive answer to the second. Since the questions differ, the answers are compatible. We show that Merleau-Ponty's objections, from the crag, from sedimented habit and from class, establish limits on what an agent can bring about and on how likely she is to change her projects, but not that facts determine meaning. The absoluteness so understood is neutral on determinism.
\end{abstract}

\section{The problem}

In Part~IV of \emph{Being and Nothingness}, Sartre argues that human reality is free in a radical sense: it is not first something and then free, but its being is its freedom, the perpetual detachment from what it is by which it projects itself towards what it is not \citep[Part~IV, ch.~1]{Sartre1943}. There are no limits to freedom except freedom itself; we are not free to cease being free. In the same chapter, under the heading \enquote{Freedom and Facticity: The Situation}, Sartre examines what appear to be limits: my place, my past, my environment, my fellow human beings and my death. A crag is too steep for me to climb; my past is what it is and cannot be changed; I was born a worker, a Frenchman, in this body. Sartre's reply is that each of these is a given that freedom did not create but that has the meaning it has only in the light of a project. The crag is an obstacle only to one who wants to climb it; to one who wants to paint the landscape it is a feature of the view. Its \enquote{coefficient of adversity} is real but relative to the project \citep[Part~IV, ch.~1]{Sartre1943}.

Merleau-Ponty objected in the final chapter of the \emph{Phenomenology of Perception} \citep{MerleauPonty1945}. If freedom is equal in all our acts, there is no free action, since free action stands out against a background of life that is not free; if every situation is equally the creation of freedom, freedom is nowhere. The crag is not big only for one who wants to climb it; it is big relative to my body, prior to any explicit project. My habits, my class and my history are sedimented in me and give my freedom a field, with lines of probability, so that some changes are not impossible but unlikely. Freedom, Merleau-Ponty concludes, is real but bounded and gradual. Beauvoir pressed a related point from the side of oppression: an enslaved woman in a harem is not free in the same way as a philosopher, since her situation conceals from her the very possibility of revolt \citep{Beauvoir1947}. Sartre himself came to qualify his early view \citep{Sartre1960}. Recent commentary divides between those who think the early view is incoherent and was rightly abandoned and those who think it was misread \citep{Kruks1990,Webber2009}. We defend a reading on which the early view is coherent.

\section{Framework}

\begin{definition}[Facticity]\label{d:fact}
The \emph{facticity} of an agent at a time is the totality of facts about her and her circumstances that hold independently of her present project: her body, place, past, environment, the existence and conduct of others, and her mortality.
\end{definition}

\begin{definition}[Meaning of a fact]\label{d:mean}
The \emph{meaning} of a fact for an agent is its practical significance in her situation: whether it is an obstacle, a means, an opportunity, a danger, an irrelevance, and in what degree.
\end{definition}

\begin{definition}[Ontological and practical freedom]\label{d:free}
An agent is \emph{ontologically free} if no fact of her facticity determines, by itself, its meaning for her. An agent's \emph{practical freedom} at a time is the range of outcomes she can bring about at that time.
\end{definition}

\begin{principle}[Meaning is project-relative]\label{p:rel}
The meaning of a fact for an agent is a function of the fact together with the agent's projects: if the projects differ, the meaning of the same fact can differ.
\end{principle}

\begin{principle}[Projects are not fixed by facts]\label{p:proj}
For any agent and any fact of her facticity, that fact does not by itself fix which project she has; her projects are fixed in the course of her existence by the ends she adopts and pursues.
\end{principle}

Principle~\ref{p:rel} is Sartre's analysis of the coefficient of adversity. Principle~\ref{p:proj} is the core of EXI-09's thesis that for human beings existence precedes essence, read, as EXI-09-a reads it, as the thesis that a person's defining project is fixed by no determination prior to her existence but by the agent who lives it. The principle does not say that projects are uncaused, or that they arise from nothing; it says that no single fact of facticity settles them, since what any given fact contributes to them depends on how it is taken up.

\section{The argument}

\begin{lemma}[Absoluteness]\label{l:abs}
Every agent is ontologically free, and ontological freedom does not come in degrees.
\end{lemma}

\begin{proof}
Let $F$ be a fact of an agent's facticity. By Principle~\ref{p:rel}, the meaning of $F$ for her is fixed by $F$ together with her projects, and could differ if her projects differed. By Principle~\ref{p:proj}, $F$ does not by itself fix her projects. So $F$ does not by itself fix its meaning for her, and since $F$ was arbitrary, the agent is ontologically free (Definition~\ref{d:free}). The property is universal: it holds of every fact for every agent who has projects. Hence it admits no degrees. An agent cannot be more or less such that facts do not determine their own meaning for her, any more than a number can be more or less even.
\end{proof}

\begin{lemma}[Facticity]\label{l:fact}
Practical freedom is limited by facticity and comes in degrees.
\end{lemma}

\begin{proof}
What an agent can bring about depends on her body, resources, position and the conduct of others, which are facts of her facticity (Definition~\ref{d:fact}). Agents differ in these, so they differ in the range of outcomes they can bring about; the climber with a broken leg can bring about less than the healthy one. So practical freedom is limited and graded.
\end{proof}

\begin{theorem}[No tension]\label{t:notension}
Ontological freedom is absolute, practical freedom is limited by facticity, and the two claims are consistent.
\end{theorem}

\begin{proof}
The first two claims are Lemmas~\ref{l:abs} and~\ref{l:fact}. They are consistent because they concern different relations. Ontological freedom concerns the relation between facts and their meaning: that facts do not by themselves determine it. Practical freedom concerns the relation between facts and outcomes: which outcomes an agent can bring about given the facts. A fact can limit outcomes without fixing its meaning: the crag limits which routes I can take, but whether that limitation is an obstacle, a challenge or an irrelevance depends on whether I am climbing, training or painting. Conversely, that meaning is project-relative does not enlarge the range of outcomes. Hence no case of facticity limiting practical freedom is a case of a fact fixing its own meaning, and the claims cannot conflict.
\end{proof}

The theorem is what Sartre calls the paradox of freedom: there is freedom only in a situation, and there is a situation only through freedom \citep[Part~IV, ch.~1]{Sartre1943}. Once the two relations are distinguished the paradox is not a contradiction but a mutual dependence. Facticity gives freedom something to take up; freedom gives facticity its meaning.

The absoluteness so defined is an ontological claim about the relation between facts and meaning. It is not a claim about the causal history of choices. Whether choices are determined by prior events is a further question, which ACT-01 answers by arguing that free will is compatible both with determinism and with indeterminism. Theorem~\ref{t:notension} is neutral on it: if determinism is true, the projects an agent adopts may have sufficient causes, but the meanings of her facts are still fixed by those projects and not by the facts alone. The view therefore makes no libertarian commitment. Nor does Principle~\ref{p:proj} require that the agent create herself from nothing; ACT-08 argues that self-formation need not start from nothing, and Principle~\ref{p:proj} requires only that it not be fixed by any one given.

\section{Merleau-Ponty's objections}

\begin{objection}[Freedom everywhere is nowhere]\label{o:nowhere}
If freedom is equally present in every act, there is no difference between a free act and an unfree one, and the concept loses its point \citep{MerleauPonty1945}.
\end{objection}

\begin{reply}
Ontological freedom is equally present in every act, as Lemma~\ref{l:abs} says. The distinction between more and less free acts is drawn by practical freedom and by the reflective appropriation of one's projects, which come in degrees (Lemma~\ref{l:fact}). The concept of ontological freedom does not have the point of discriminating among acts; its point is to characterize the kind of being for which acts have meanings at all. A concept that applies to every case of its domain is not thereby empty, as the concept of having a shape applies to every body without being idle.
\end{reply}

\begin{objection}[The crag is big for my body]\label{o:crag}
Before any project, the crag appears as big relative to my body; its size is a meaning that the project does not confer \citep{MerleauPonty1945}.
\end{objection}

\begin{reply}
We grant that the crag appears as big relative to the body's capacities before any explicit project, and that this is a meaning in a broad sense. But it is a meaning relative to the body's general project of moving in the world, which Merleau-Ponty himself calls an anonymous, general existence. That is still a project, in the sense of Principle~\ref{p:rel}: the crag's size matters as it does because the body is a being that moves, climbs and reaches. Merleau-Ponty's point shows that some meanings are fixed by projects one has not chosen explicitly, which is a point about the generality of the body's projects, not a case of a fact fixing its own meaning. PHE-08 argues that the lived body is the zero-point of orientation, from which size and distance are given; on that account size is relative to the body, not to the crag alone.
\end{reply}

\begin{objection}[Sedimentation]\label{o:sed}
Habits and the past are sedimented in us, so that some changes of project, though not impossible, are improbable. Freedom therefore has a field with lines of force, and is not absolute \citep{MerleauPonty1945}.
\end{objection}

\begin{reply}
The objection concerns the probability of changing one's projects, which is a matter of practical freedom with respect to oneself. It is true and important: an inferiority complex of twenty years' standing is unlikely to be abandoned in a day. But that improbability is itself a fact of facticity, and its meaning for the agent, whether it is a fate, a challenge or a reason for patience, depends on her projects. Sedimentation thus limits what she is likely to bring about without determining what it means. Merleau-Ponty's corrective is correct as an account of practical freedom, and Sartre's early account is weak on it; but the two accounts are answering different questions.
\end{reply}

\begin{objection}[Oppression]\label{o:opp}
The oppressed woman in Beauvoir's example cannot even conceive revolt. To call her free absolutely is an insult.
\end{objection}

\begin{reply}
Beauvoir's point is that oppression can narrow not only what one can do but what possibilities appear to one \citep{Beauvoir1947}. That narrowing is a limitation of practical freedom, including the practical freedom to form projects, and it is a wrong done to her. It does not show that the facts of her situation fix their own meaning; it shows that others have arranged her situation so that only some meanings are available. The ontological claim is not a claim about her opportunities, and it carries no suggestion that her situation is her responsibility in the way that a free choice among open options would be. The limits of bad faith are taken up in PHE-10-b.
\end{reply}

\section{What the result does not establish}

The result does not establish that agents are responsible for everything that befalls them, which Sartre sometimes asserts; responsibility depends on practical freedom as well. It does not establish libertarian free will or refute determinism. And it does not establish that all projects are equally available or equally good; it says only that whatever projects an agent has, the facts of her situation do not by themselves fix their meaning.

\section{Conclusion}

Freedom is ontologically absolute: no fact of an agent's facticity fixes by itself what it means for her, and this holds of every agent without degrees. Freedom is always exercised within facticity: what an agent can bring about, and how likely she is to change her projects, are limited by her body, past, place and others. The two claims concern different relations, between facts and meaning and between facts and outcomes, and so are not in tension. Merleau-Ponty and Beauvoir were right about practical freedom and its social conditions; they did not show that the absoluteness of ontological freedom is false.

\begin{thebibliography}{99}\small
\bibitem[Beauvoir(1947)]{Beauvoir1947} Beauvoir, S. de (1947). \emph{Pour une morale de l'ambigu\"it\'e}. Gallimard. Trans. B. Frechtman, \emph{The Ethics of Ambiguity}, Philosophical Library, 1948.
\bibitem[Kruks(1990)]{Kruks1990} Kruks, S. (1990). \emph{Situation and Human Existence: Freedom, Subjectivity and Society}. Unwin Hyman.
\bibitem[Merleau-Ponty(1945)]{MerleauPonty1945} Merleau-Ponty, M. (1945). \emph{Ph\'enom\'enologie de la perception}. Gallimard. Trans. D.~A. Landes, \emph{Phenomenology of Perception}, Routledge, 2012.
\bibitem[Sartre(1943)]{Sartre1943} Sartre, J.-P. (1943). \emph{L'\^etre et le n\'eant}. Gallimard. Trans. H.~E. Barnes, \emph{Being and Nothingness}, Philosophical Library, 1956.
\bibitem[Sartre(1960)]{Sartre1960} Sartre, J.-P. (1960). \emph{Critique de la raison dialectique}. Gallimard.
\bibitem[Webber(2009)]{Webber2009} Webber, J. (2009). \emph{The Existentialism of Jean-Paul Sartre}. Routledge.
\end{thebibliography}

\end{document}
